The thesis shows that pit stop decision support can be studied from historical Formula 1 data only if the problem is made 
\textbf{temporal, contractual, and comparable}. Historical records must be replayed or reconstructed in event order; the target must state exactly 
what a correct positive recommendation means; and systems must be evaluated on the same truth universe. Without these constraints, the comparison 
between a rule engine, a Batch learner, and an online learner would mix modelling quality with leakage, denominator differences, or post-race information.

The first research question receives a positive answer. Near-term pit timing, represented by \textbf{\texttt{pit\_any\_h2}}, is a \textbf{learnable task}. 
The learning systems improve over the deterministic baseline, with Batch Percent providing the largest event coverage and MOA offering the 
best precision-oriented operating points. The result shows that pit stops can't be predicted perfectly, but race state patterns contain 
enough information to support better than rule-baseline timing recommendations under a live like chronology.

The second research question receives a more cautious answer. Successful pit recommendation, represented by \textbf{\texttt{pit\_success\_h2}}, 
is \textbf{much harder}. MOA Percent obtains the cleanest selected strict operating point, but coverage remains low. Relaxed thresholds can recover
more events, at the cost of more false positives. This confirms that successful pit prediction is harder than pit timing: 
it depends on later race development, rival reactions, traffic, tyre performance, and the imperfect heuristic used to classify pit outcomes.

The main limitations follow from the experimental nature of the work. The replay is controlled and reproducible, but it is not a production live 
feed with all latency, correction, missing-data, and failure modes. The PitStopEvaluator is a useful approximation of outcome quality, not yet a full 
strategic utility model. Success labels are sparse, so race-level diagnostics can be unstable. Thresholds strongly influence the apparent behaviour of 
each system. Finally, only a limited set of learners is evaluated: XGBoost for Batch learning and OzaBoostADWIN for MOA online learning.

Future work should therefore move in four directions. First, the pipeline should be connected to a live ingestion setting, so that delay and 
missing-data behaviour can be evaluated directly. Second, the truth model should be refined through richer outcome definitions, probabilistic labels, 
or expert validation. Third, the learning problem could move from row-level classification toward event-level ranking, time-to-pit modelling, or 
sequence-aware policies under a fixed false-positive budget. Fourth, online learning should be expanded with stronger calibration and faithful 
online explainability.

Overall, the contribution of the thesis is an \textbf{executable and auditable framework} for asking a precise question: when a system says ``pit now'', 
what contract is it satisfying, what information was available, and how should that call be counted? The answer is strongest for pit timing and still 
preliminary for successful pit recommendation, but the pipeline makes both problems measurable in a reproducible way.
